\subsection{Upper envelopes and sums of measures}

PROPOSITION 6. --- Let \((\lambda_{\alpha})_{\alpha \in A}\) be an increasing directed family of measures on T, and let \(p = \sup_{\alpha} \lambda_{\alpha}^{\bullet}\) . For the family \((\lambda_{\alpha})\) to be bounded above in \(\mathcal{M}(T)\) , it is necessary and sufficient that the encumbrance \(p\) be locally bounded. The family \((\lambda_{\alpha})\) then admits a supremum \(\lambda\) in \(\mathcal{M}(T)\) , and \(\lambda^{\bullet} = p\) . For every compact set K, the measure \(\lambda_{K}\) is the supremum of the measures \((\lambda_{\alpha})_{K}\) in \(\mathcal{M}(K)\) .

If the family \((\lambda_{\alpha})\) is bounded above in \(\mathcal{M}(\mathrm{T})\) , then p is obviously locally bounded. Conversely, let us assume p to be locally bounded, and let us show that it satisfies the conditions 1) and 2) of Prop. 2 b) of No.~2. For 2), this results from the following equalities:

\[
p (f) = \sup _ {\alpha} \lambda_ {\alpha} ^ {\bullet} (f) = \sup _ {\alpha} \sup _ {K} \lambda_ {\alpha} ^ {\bullet} (f \varphi_ {K}) = \sup _ {K} \sup _ {\alpha} \lambda_ {\alpha} ^ {\bullet} (f \varphi_ {K}) = \sup _ {K} p (f \varphi_ {K}).
\]

On the other hand, let K be a compact set; the encumbrance \(p_{K}\) is equal to the upper envelope of the encumbrances \((\lambda_{\alpha}^{\bullet})_{K}\) and it is bounded since p is locally bounded. The measures \((\lambda_{\alpha})_{K}\) therefore admit a supremum \(\lambda_{K}\) in \(\mathcal{M}(K)\) , and \(\lambda_{K}^{\bullet}=p_{K}\) (Ch. V, §1, No.~4, Prop. 11). The condition 1) of Prop. 2 b) of No.~2 is thus satisfied, therefore there exists a measure \(\lambda\) on T such that \(\lambda^{\bullet}=p\) ; it is clear that \(\lambda\) is the supremum of the measures \(\lambda_{\alpha}\) .

DEFINITION 9. --- Let \((\mu_i)_{i \in I}\) be a family of measures on T. Let A be the set of finite subsets of I; for every \(\alpha \in A\) let \(\lambda_\alpha = \sum_{i \in \alpha} \mu_i\) . If the family \((\lambda_\alpha)\) admits a supremum \(\mu\) in \(\mathcal{M}(T)\) , the family \((\mu_i)\) is said to be summable, \(\mu\) is called the sum of the family \((\mu_i)\) , and one writes \(\mu = \sum_{i \in I} \mu_i\) .

This definition extends the definition of Ch. V, §2, No.~1.

PROPOSITION 7. --- For the family \((\mu_i)_{i \in I}\) to be summable, with sum \(\mu\) , it is necessary and sufficient that the encumbrance \(p = \sum_{i \in I} \mu_i^\bullet\) be locally bounded, in which case \(p = \mu^\bullet\) . For every compact subset K of T, the family \(\left((\mu_i)_K\right)_{i \in I}\) is then summable in \(\mathcal{M}(K)\) , and \(\mu_K = \sum_{i \in I} (\mu_i)_K\) .

With notations as in Def. 9, \(\lambda_{\alpha}^{\bullet} = \sum_{i\in \alpha}\mu_{i}^{\bullet}\) for every finite subset \(\alpha\) of I (No.~2, Remark 1). The statement is then an immediate consequence of Prop. 6.

The relation \(\mu_{\mathrm{K}} = \sum_{i\in \mathrm{I}}(\mu_i)_{\mathrm{K}}\) and Prop. 2 of Ch. V, §2, No.~2 yields the following result:

PROPOSITION 8. --- Let \(\mu\) be the sum of a summable family \((\mu_i)_{i \in I}\) of measures on T. In order that a mapping \(f\) of T into a topological space F (Hausdorff or not) be \(\mu\) -measurable, it is necessary and sufficient that \(f\) be \(\mu_i\) -measurable for every \(i \in I\) .

